\section{Rings and Modules}

\subsection{Rings and ideals}

We recall some basic notions regarding rings.

\bdefn

    A {\emphcolor ring} is a set $R$ equipped with two operations $+$ and $\cdot$, making it into a commutative group and a monoid respectively.
    Furthermore, they satisfy the distributivity properties:
    $$ a\cdot(b+c) = ab + ac,\qquad (a+b)\cdot c = ac + bc $$
    for all $a,b,c\in R$.
    The additive and multiplicative identities are denoted $0$ and $1$ respectively.

\edefn

\bnote

    In this course all rings are assumed to be {\emphcolor commutative}: $\cdot$ forms a commutative monoid.

\enote

\bdefn

    A {\emphcolor ring morphism} is a function between two rings, $f\colon R\to S$, which preserves structure: $f(x+y)=f(x)+f(y)$, $f(xy)=f(x)f(y)$,
    and $f(1)=1$.
    A {\emphcolor ring isomorphism} is a ring morphism with an inverse (equivalently a bijective ring morphism).

    A {\emphcolor subring} of a ring $R$ is a subset $S\subseteq R$ equipped with a ring structure making the inclusion $S\to R$ a ring morphism.

    An {\emphcolor ideal} of a ring $R$ is a nonempty subset $I\subseteq R$ closed under addition and multiplication by $R$: for $r\in R$ and $x\in I$,
    $rx\in I$.
    For an ideal $I\subseteq R$, we define its {\emphcolor quotient ring} $R/I$ to be the set of all cosets $[r]=r+I$ with ring structure given by
    $$ 0 = [0] = I,\qquad 1 = [1],\qquad [x] + [y] = [x+y],\qquad [x]\cdot[y] = [xy] , $$
    equivalently the unique structure making the projection map $\pi\colon R\to R/I,r\mapsto[r]$ a ring morphism.

\edefn

\bprop

    Let $f\colon R\to S$ be a ring morphism, then
    \benum
        \item For every ideal $J\leq S$, $f^{-1}(J)\leq R$ is an ideal.
        \item In particular $f$'s {\emphcolor kernel}, $\ker f=f^{-1}0$, is an ideal of $R$.
        And $f$'s {\emphcolor image}, $\im f=f(R)$ is a subring of $S$.
        $f$ then quotients to give a ring isomorphism $R/\ker f\cong\im f$.
        \item For an ideal $I\leq R$, consider the projection $\pi\colon R\to R/I$.
        Then $J\mapsto\bar J=\pi(J)$ and $\bar J\mapsto\pi^{-1}\bar J$ gives a one-to-one correspondence between ideals of $R$ containing $I$
        and ideals of $R/I$.
    \eenum

\eprop

\bdefn

    For a ring $R$ and a subset $S\subseteq R$, we denote the smallest ideal containing $S$ by $(S)$.
    This exists as the arbitrary intersection of ideals is an ideal.
    When $S=\set x$ is a singleton, we denote $(\set x)=(x)$; such an ideal is called a {\emphcolor principal ideal}.

\edefn

\bprop

    The ideal generated by $S$, $(S)$, is the set of all finite $R$-combinations of elements from $S$:
    $$ (S) = \set{\sum_{s\in S}r_ss}[r_s\in R,\hbox{ all but finitely many being zero}] . $$
    In particular, $(x)=Rx$ is the set of all multiples of $x$.

\eprop

\bproof

    It is not hard to see that this is an ideal, and clearly every ideal containing $S$ must contain it.
    \qed

\eproof

\bdefn

    Let $R$ be a ring, an element $x\in R$ is
    \benum
        \item a {\emphcolor zero divisor} if there is a $y\in R$ such that $xy=0$;
        \item a {\emphcolor unit} if there is a $y\in R$ such that $xy=1$;
        \item {\emphcolor nilpotent} if $x^n=1$ for some $n>0$.
    \eenum
    The set of units of $R$ is denoted $R^\times$ and forms a group.

\edefn

\bdefn

    A ring $R\neq0$ is
    \benum
        \item {\emphcolor reduced} if every $x\neq0$ is not nilpotent;
        \item an {\emphcolor integral domain} if every $x\neq0$ is not a zero divisor;
        \item a {\emphcolor field} if every $x\neq0$ is a unit.
    \eenum

\edefn

It is easy to see that a field is an integral domain is reduced.

\bprop

    For a ring $R$, the following are equivalent:
    \benum
        \item $R$ is a field,
        \item $R$ has no non-trivial ideals,
        \item Every ring morphism out of $R$ to a non-zero ring is injective.
    \eenum

\eprop

\bproof

    If $R$ is a field, then every nonzero element is a unit and thus the ideal generating it must be all of $R$.
    Conversely if $R$ has no non-trivial ideals, then for $x\neq0$ we must have that $1\in(x)$ and thus there is a $y\in R$ such that $xy=1$ so it is a unit.

    If $R$ has no non-trivial ideals then for every $f\colon R\to S\neq0$ we must have $\ker f=0$, as it is an ideal not containing $1$.
    Thus $f$ is injective.
    Conversely if $I\neq R$ is an ideal of $R$, then $\pi\colon R\to R/I$ is injective, but its kernel is $I$, so $I=0$.
    \qed

\eproof

\bprop

    If $\set{I_\alpha}_\alpha$ is a collection of ideals from $R$, then
    \benum
        \item $\bigcap_\alpha I_\alpha$ is an ideal, the largest ideal contained in each $I_\alpha$.
        \item $\sum_\alpha I_\alpha=\set{\sum_\alpha x_\alpha}[x_\alpha\in I_\alpha\hbox{ all but finitely many are zero}]$ is an ideal, the smallest ideal containing each $I_\alpha$.
    \eenum

\eprop

\bproof

    The first point is easily verified, and the second one follows because
    $$ \sum_\alpha I_\alpha = \parens{\bigcup_\alpha I_\alpha} . \qed $$

\eproof

\bdefn

    If $I,J$ are ideals then we define their {\emphcolor product} to be the ideal generated from all products of elements from $I$ and $J$:
    $$ IJ = (\set{xy}[x\in I,y\in J]) . $$

\edefn

Since $I,J$ are ideals and are thus closed under multiplication from $R$, for $x\in I$ and $y\in J$ we have that $xy\in I\cap J$.
Thus $IJ\subseteq I\cap J$.

\bdefn

    For an ideal $I\subseteq R$, we define its {\emphcolor radical} to be
    $$ \sqrt I = \set{r\in R}[r^n\in I\hbox{ for some $n\geq0$}] . $$

\edefn

\blemm

    The radical of an ideal is an ideal.

\elemm

\bproof

    If $a^n\in I$ and $r\in R$ then $(ra)^n=r^na^n\in I$ so that $ra\in\sqrt I$.
    And if $a^n,b^m\in I$ then
    $$ (a+b)^{n+m} = \sum_{i=0}^{n+m}\binom{n+m}i a^ib^{n+m-i} \in I , $$
    since for each $i$, if $i<n$ then $n+m-i>m$, so $a+b\in\sqrt I$.
    \qed

\eproof

\subsection{Prime and maximal ideals}

\bdefn

    A proper ideal $I\subset R$ is
    \benum
        \item {\emphcolor prime} if $ab\in I$ implies $a\in I$ or $b\in I$.
        Equivalently if $a,b\notin I$ then $ab\notin I$.
        \item {\emphcolor maximal} if it is not contained in any other proper ideal.
    \eenum

\edefn

\bprop

    An ideal $I$ is
    \benum
        \item prime iff $R/I$ is an integral domain;
        \item maximal iff $R/I$ is a field.
    \eenum
    Consequently, a maximal ideal is prime.

\eprop

\bproof

    $R/I$ is an integral domain iff $[a][b]\neq[0]$ when $[a],[b]\neq[0]$.
    This means that $ab\notin I$ when $a,b\notin I$, as desired.

    $R/I$ is a field iff it has no nontrivial ideals.
    Since ideals of $R/I$ are in bijection with ideals containing $I$, this is iff $I$ is maximal.
    \qed

\eproof

In particular the zero ideal $0=(0)=\set0$ is prime iff $R\cong R/(0)$ is an integral domain.

\bcoro

    Let $I$ be an ideal, then the correspondence $J\oto\bar J$ between ideals containing $I$ and ideals of $R/I$
    preserves primality and maximality: $J$ is prime/maximal iff $\bar J$ is.

\ecoro

\bproof

    We have that
    $$ (R/I)/\bar J = (R/I)/(J/I) \cong R/J , $$
    so that $(R/I)/\bar J$ is an integral domain/field iff $R/J$ is.
    The result follows.
    \qed

\eproof

\bdefn

    \benum
        \item An element $0\neq x\in R$ is {\emphcolor irreducible} if it is not a unit, and if $x=yz$ then either $y$ or $z$ is a unit.
        \item An integral domain $R$ is a {\emphcolor principal ideal domain} (or a {\emphcolor PID}) if all of its ideals are principal.
    \eenum

\edefn

\bprop

    If $R$ is a PID and $0\neq x$, then $(x)$ is prime iff $x$ is irreducible iff $(x)$ is maximal.

\eprop

\bproof

    If $(x)$ is prime, $x$ is not a unit since it is a proper ideal.
    And if $x=yz$ we have that $y$ or $z$ is in $(x)$, say $y=xt$.
    Then $x=yz=xtz$, and so $x(1-tz)=0$, and since $R$ is an integral domain this means $tz=1$ so $z$ is a unit.

    Now if $x$ is irreducible, then suppose $(x)\subseteq I$ is a proper ideal containing $(x)$.
    Since $R$ is a PID, $I=(y)$ so that $x\in(y)$ meaning $x=yz$.
    $(y)$ is a proper ideal so $y$ is not a unit, thus $z$ is, and then $y=xz^{-1}$ so that $y\in(x)$ and thus $(x)=(y)$.
    So $(x)$ is maximal, as desired.
    \qed

\eproof

\bexam

    The ring of integers $\bZ$ is a PID.
    Its irreducible elements are precisely its prime elements, so that its prime ideals are $(0)$ and $(p)$ for $p\in\bZ$ prime.

\eexam

\bexam

    If $k$ is a field, then its ring of polynomials $k[x]$ is a PID.
    Indeed, if $0\neq I\leq k[x]$ is an ideal, then let $f\in I$ have minimal degree.
    By polynomial division we see that $I=(f)$: for $g\in I$ we have that $g=fq+r$, but since $r=g-fq\in I$ has smaller degree than $f$,
    it must be zero.

    Its prime ideals are $(0)$ and $(f)$ for an irreducible monic polynomial $f$.

\eexam

\blemm

    If $f\colon R\to S$ is a ring morphism, and $\q\leq S$ is prime, then $f^{-1}\q\leq R$ is prime.

\elemm

\bproof

    A direct proof is easy, instead consider the composition:
    $$ \bar f\colon R\xtos fS\xtos\pi S/\q . $$
    By definition $f^{-1}\q=\ker\bar f$, and so $\bar f$ induces an embedding $R/f^{-1}\q\inj S/\q$.
    Subrings of integral domains are integral domains, so $R/f^{-1}\q$ is an integral domain, thus $f^{-1}\q$ is prime.
    \qed

\eproof

Subrings of fields are not fields, so the above proof does not show that the preimage of a maximal ideal is maximal.
Indeed, this is not true: consider the unique map $f\colon\bZ\to\bQ$.
Then $0\leq\bQ$ is maximal, but $f^{-1}0=\ker f=0$ is not maximal in $\bZ$.

\bthrm

    For every ring $R$ and every proper ideal $I<R$, there is a maximal ideal $\m$ containing $I$.

\ethrm

\bproof

    Consider the set of all proper ideals containing $I$.
    This is nonm-empty, as it contains $I$.
    Furthermore, the union of every chain of proper ideals is a proper ideal; it being an ideal is easy to see, and it is proper
    since $1$ is not in any element of the chain.
    Thus by Zorn, this set has a maximal element, which must be a maximal ideal.
    \qed

\eproof

\bcoro

    Let $x\in R$, then $x$ is a unit iff $x$ is not in any maximal ideal.
    That is,
    $$ R^\times = \bigcap_\m\m^c , $$
    where the intersection is over all maximal ideals of $R$.

\ecoro

\bproof

    An element is a unit iff it is not contained in any proper ideal.
    This is clear enough: if $x$ is not a unit then $(x)$ is proper, and proper ideals don't contain units.
    Every proper ideal can be extended to a maximal ideal, so the result follows.
    \qed

\eproof

\bdefn

    Let $R$ be a ring.
    Then define its {\emphcolor nilradical} to be $\nil(R)=\sqrt0$, i.e.\ the ideal of all nilpotent elements.

\edefn

Clearly $R$ is reduced iff $\nil(R)=0$.

\bprop

    Let $I\leq R$ be an ideal, and let $\pi\colon R\to R/I$ be the projection.
    Then $\sqrt I=\pi^{-1}\nil(R/I)$.

\eprop

\bproof

    We note that if $x^n\in I$ then $\pi(x)^n\in I$ so that $\pi(x)\in\nil(R/I)$ as desired.
    Conversely, if $\pi(x)\in\nil(R/I)$ then $\pi(x^n)=I$ so that $x^n\in I$ as desired.
    \qed

\eproof

\bprop

    For a ring $R$,
    \benum
        \item $R/\nil(R)$ is reduced;
        \item The nilradical is the intersection of all prime ideals:
        $$ \nil(R) = \bigcap_\p\p . $$
    \eenum

\eprop

\bproof

    We noted above that $\nil R=\sqrt{\nil R}=\pi^{-1}\nil(R/\nil(R))$.
    But we have also that $\nil R=\pi^{-1}0$, so $\pi^{-1}0=\pi^{-1}\nil(R/\nil(R))$ and since $\pi$ is surjective this means $\nil(R/\nil(R))=0$
    so that $R/\nil(R)$ is reduced.

    Clearly if $x\in\nil(R)$ then $x\in\p$ for every prime $\p$ since $x^n=0\in\p$.
    Now if $x\notin\nil(R)$ then we claim there exists a prime ideal $\p$ such that $x\notin\p$.
    We prove this using Zorn.
    Let $T=\set{1,x,x^2,\dots}$ and let $S$ be the poset of all ideals $I$ such that $I\cap T=\emptyset$; note that all ideals in $S$ are proper.
    As explained before the union of a chain of ideals is an ideal, and clearly if all of the ideals in the union don't intersect $T$, neither
    does the union.
    Thus by Zorn, $S$ has a maximal element $\p$.

    We claim that $\p$ is prime.
    Indeed, suppose that $ab\in\p$ but $a,b\notin\p$.
    Then by maximality we have that $x^n\in(a,\p)$ and $x^m\in(b,\p)$ for some $n,m\geq0$.
    But then $x^{n+m}\in(a,\p)(b,\p)\subseteq\p$, a contradiction.

    Thus there is a prime ideal $\p$ such that $\p\cap T=\emptyset$, in particular $x\notin\p$.
    So if $x\notin\nil(R)$ then $x\notin\bigcap_\p\p$, as desired.
    \qed

\eproof

\bcoro

    Let $I$ be an ideal of $R$, then
    $$ \sqrt I = \bigcap_{I\leq\p}\p , $$
    where the intersection is over all prime ideals containing $I$.

\ecoro

\bproof

    We showed above that $\sqrt I=\pi^{-1}\nil(R/I)$.
    By above $\nil(R/I)$ is the intersection over all prime ideals of $R/I$, which are all prime ideals of $R$ containing $I$:
    $$ \sqrt I = \pi^{-1}\bigcap_{I\leq\p}\bar\p = \bigcap_{I\leq\p}\pi^{-1}\bar\p = \bigcap_{I\leq\p}\p . \qed $$

\eproof

\bdefn

    Let $R$ be a ring, then define its {\emphcolor Jacobson radical} to be
    $$ J(R) = \bigcap_\m\m , $$
    the intersection of all maximal ideals.

\edefn

\blemm

    Let $\m\leq R$ be maximal.
    Then $x\notin\m$ iff there is a $y\in R$ such that $1-xy\in\m$.

\elemm

\bproof

    $x\notin\m$ iff $\bar x\in R/\m$ is nonzero, or equivalently is a unit.
    This is iff there is a $y\in R$ such that $\bar x\bar y=\bar 1$.
    This just means $1-xy\in\m$, as desired.
    \qed

\eproof

\bcoro

    For a ring $R$, its Jacobson radical $J(R)$ is the set of all elements $x\in R$ such that $1-yx\in R^\times$ for all $y\in R$.

\ecoro

\bproof

    Now if $x\notin J(R)$ then $x\notin\m$ for some $\m$, and then $1-yx\in\m$ for some $y$ and thus is not a unit.
    And if $1-yx$ is not a unit, then $1-yx\in\m$ for some maximal $\m$ and thus $x\notin\m$ so that $x\notin J(R)$.
    \qed

\eproof

\bdefn

    A ring $R$ is {\emphcolor local} if it has a unique maximal ideal $\m$.
    The field $k=R/\m$ is called its {\emphcolor residue field}.
    By abuse of notation we write $(R,\m)$ for a local field.

\edefn

\bprop

    \benum
        \item If $(R,\m)$ is a local ring then $R-\m\subseteq R^\times$.
        \item Conversely if $I<R$ is a proper ideal with $R-I\subseteq R^\times$ then $(R,I)$ is a local ring.
        \item If $\m\leq R$ is maximal such that $1+\m\subseteq R^\times$ then $(R,\m)$ is local.
    \eenum

\eprop

\bproof

    If $x\notin\m$ then $(x)$ is not contained in $\m$, and since $\m$ is the unique maximal ideal this means that $(x)=R$ so that $x$ is a unit.
    Also we noted that $R^\times$ is the intersection of all complements of maximal ideals, so this follows from that as well.

    Conversely, we have that
    $$ R-I \subseteq R^\times \bigcap_\m\m^c \implies I \supseteq \bigcup_\m\m , $$
    which means that $I$ contains all maximal ideals, which can only happen if it is the unique maximal ideal.

    We want to show that $R-\m\subseteq R^\times$.
    Indeed, if $x\notin\m$ then there is a $y$ such that $1-yx\in\m$ so that $yx\in1+\m$ and is thus a unit.
    So $x$ is a unit, as desired.
    \qed

\eproof

\bexam

    For $p\in\bZ$ prime, define
    $$ \bZ_{(p)} = \set{\frac ab}[a,b\in\bZ,b\notin(p)] \subseteq \bQ ,\qquad I = \set{\frac ab\in\bZ_{(p)}}[a\in(p)] . $$
    Then $(\bZ_{(p)},I)$ is local.
    Indeed, if $\frac ab\notin I$ then it is invertible (its inverse is $\frac ba$).
    We can generalize this, as we will in the next section.

\eexam

\subsection{Modules}

\bdefn

    Let $A$ be a ring, then a {\emphcolor module} over $A$ is an Abelian group $M$ along with a ring morphism
    $\phi\colon A\to\endo(M)$.
    We define {\emphcolor multiplication} to be a function $\mu\colon A\times M\to M$, written by juxtaposition, by
    $$ \mu(a,m) = \phi(a)m . $$

\edefn

An $A$-module can also be defined via multiplication: require that multiplication distributes both ways, and $1\in A$ is an identity:
$$ (a+b)m = am + bm,\qquad a(m+n) = am + an,\qquad a(bm) = (ab)m,\qquad 1m = m . $$

\bdefn

    For two $A$-modules $M$ and $N$, an {\emphcolor $A$-module morphism} is a map $f\colon M\to N$ which is an Abelian group morphism and also
    multiplicative:
    $$ f(am) = af(m) $$
    for all $a\in A,m\in M$.

\edefn

\bnote

    I'm going to skip the rest of the standard definitions, for Chrissake.

\enote

\bdefn

    If $M$ is an $A$-module, we define $\ann_A(M)\leq A$ the {\emphcolor annihilator} of $M$ by
    $$ \ann_A(M) = \set{a\in A}[aM=0] = \ker\phi , $$
    where $\phi$ is the module map $\phi\colon A\to\endo(M)$.
    This is then an ideal.
    For $S\subseteq M$ a subset we define $\ann_A(S)$ similarly.

    A module $M$ is {\emphcolor faithful} if $\ann_A(M)=0$.

\edefn

$\ann_A(S)$ is equal to $\ann_A(\gen S)$, the annihilator of the submodule generated by $S$ (all $R$-combinations of elements from $S$),
and is thus also an ideal.

\bdefn

    For a ring $A$, an {\emphcolor $A$-algebra} is a ring $B$ equipped with a ring morphism $\phi\colon A\to B$.

\edefn

Note that $B$ can be embedded into $\endo(B)$ mapping $b\in B$ to multiplication by $b$.
Thus an $A$-algebra is an $A$-module with multiplication given by $a\cdot b=\phi(a)b$.

\bdefn

    For two $A$-algebras $B,C$, an {\emphcolor $A$-algebra morphism} is a map $f\colon B\to C$ which is both a ring morphism and an
    $A$-module morphism.

\edefn

If $\phi\colon A\to B,\psi\colon A\to C$ are the algebra maps, then $f\colon B\to C$ is an algebra morphism iff it is a ring morphism and
$\psi=f\circ\phi$:
$$ \psi(a) = a1_C = af(1_B) = f(a1_B) = f\circ\phi(a) . $$

\bdefn

    For an $A$-module $M$ and a submodule $N$, its quotient module $M/N$ is the set of cosets $m+N$ with the usual $M$-module structure.
    The usual isomorphism results hold.

\edefn

\bprop

    Let $I\leq A$ be a module, then every $A/I$-module is an $A$-module with $\ann_A(M)\supseteq I$ and vice-versa.
    In particular an $A$-module $M$ can be viewed as a faithful $A/\ann_A(M)$-module.

\eprop

\bproof

    Consider the projection $\pi\colon A\to A/I$.
    Composing with $\phi\colon A/I\to\endo(M)$ gives an $A$-module structure to an $A/I$-module and $\ker(\phi\circ\pi)\supseteq\ker\phi=I$.
    Conversely if $\phi\colon A\to\endo(M)$ has kernel containing $I$ we can quotient it by $I$.
    \qed

\eproof

\bdefn

    For a collection of $A$-modules $\set{M_\alpha}_\alpha$ we define their {\emphcolor direct product} and {\emphcolor direct sum}
    by
    $$ \prod_\alpha M_\alpha = \set{(m_\alpha)}[m_\alpha\in M_\alpha],\qquad
    \bigoplus_\alpha M_\alpha = \set{(m_\alpha)\in\prod_\alpha M_\alpha}[\hbox{All but finitely many of $m_\alpha$ are zero}] . $$
    $\prod_\alpha M_\alpha$'s module structure is given by pointwise operations, and $\bigoplus_\alpha M_\alpha$ is viewed as a submodule.
    These form the products and coproducts in the category of $A$-modules.

\edefn

\bdefn

    An $A$-module $M$ is {\emphcolor free} if
    $$ M \cong A^{\oplus S} = \bigoplus_{\alpha\in S}A $$
    for some set $S$.
    Note that $\set{e_\alpha}_{\alpha\in S}$ (the usual standard basis vectors) forms a basis of $A^{\oplus S}$ in the sense that every
    element of $A^{\oplus S}$ can be written as a unique combination of standard basis vectors.
    A collection of elements $\set{m_\alpha}_{\alpha\in S}$ is a {\emphcolor basis} of $M$ if mapping $m_\alpha\mapsto e_\alpha$ extends
    to an isomorphism $M\cong A^{\oplus S}$.

\edefn

A basis of $M$ is equivalently a set of elements which generate $M$ and are also linearly independent in the usual way.

A free module is called such because they are free in the usual sense with respect to the obvious forgetful functor $\Mod_A\to\Set$.
That is, $\hom_A(A^{\oplus S},B)\cong\hom_\Set(S,B)$.

\bprop

    If $A^{\oplus S}\cong A^{\oplus S'}$ then $S$ and $S'$ have the same cardinality.

\eprop

\bproof

    Let $\m\leq A$ be maximal, then the isomorphism $A^{\oplus S}\cong A^{\oplus S'}$ induces an isomorphism
    $(A/\m)^{\oplus S}\cong(A/\m)^{\oplus S'}$.
    $A/\m$ is a field, and so this is an isomorphism of $A/\m$-modules (vector spaces) of dimensions $\abs S,\abs{S'}$.
    Since the dimension is invariant, their cardinalities must be equal.
    \qed

\eproof

\bdefn

    If $M$ is a free $A$-module, then its {\emphcolor rank} is the cardinality of $S$ where $M\cong A^{\oplus S}$.

\edefn

\bdefn

    An $A$-module $M$ is {\emphcolor finitely generated} ({\emphcolor fg}) if there is a finite set $S\subseteq M$ such that $M=\gen S$.

\edefn

\bprop

    An $A$-module is finitely generated iff it is a quotient of $A^n$ for some $n$ (up to isomorphism).

\eprop

\bproof

    If $\phi\colon A^n\to M$ is surjective, then $\phi(e_i)$ generate $M$ clearly and so it is finitely-generated.
    Thus if $M\cong A^n/N$ is a quotient of $A^n$ then the projection $\pi\colon A^n\to A^n/N\cong M$ is surjective and thus $M$
    is finitely-generated.
    Conversely let $m_i$ generate $M$, and define $\phi\colon A^n\to M$ by $\phi(e_i)=m_i$ which extends by linearity to all of $A^n$ since
    it is free.
    This is surjective and then $A^n/\ker\phi\cong M$ as desired.
    \qed

\eproof

\bprop

    Let $M$ be a fg $A$-module, and $I\leq A$ an ideal such that $M=IM$.
    Then there is an $x\in I$ such that $(1-x)M=0$ ($xm=m$ for all $m\in M$).

\eprop

\bproof

    Let $m_1,\dots,m_n$ generate $M$, we induct on $n$.
    For $n=0$ we have $M=0$ so this is clear.
    Consider $N=M/(m_n)$, then $IN=(IM+(m_n))/(m_n)=M/(m_n)=N$ and $N$ is generated by $n-1$ elements, so by induction there is an $x\in I$
    such that $(1-x)N=0$.
    Meaning $(1-x)M\subseteq(m_n)$, and since $m_n\in M=IM$ we have
    $$ (1-x)m_n \in (1-x)IM = I(1-x)M \subseteq I(m_n), $$
    so $(1-x)m_n=ym_n$ for some $y\in I$ and thus $(1-x-y)m_n=0$.
    Now $(1-x-y)(1-x)\in 1+I$ and satisfies
    $$ (1-x-y)(1-x)M \subseteq (1-x-y)(m_n) = 0 $$
    as desired.
    \qed

\eproof

\bcoro[title=Nakayama's lemma]

    Let $M$ be a fg $A$-module, such that $J(A)M=M$, then $M=0$.

\ecoro

\bproof

    By above there is an $x\in J(A)$ such that $(1-x)M=0$.
    Since $x\in J(A)$, $1-x\in A^\times$ so that $(1-x)M=M=0$ as desired.
    \qed

\eproof

\bcoro

    Let $M$ be a fg $A$-module and $N\subseteq M$ a submodule such that $M=IM+N$ for some $I\subseteq J(A)$, then $M=N$.

\ecoro

\bproof

    Consider $M/N$, then $I(M/N)=(IM+N)/N=M/N$, so by above $M/N=0$ meaning $M=N$.
    \qed

\eproof

\bcoro

    Let $(A,\m)$ be local, $M$ a fg $A$-module.
    Then $x_1,\dots,x_n\in M$ generate $M$ iff $\bar x_1,\dots,\bar x_n\in M/\m M$ generate $M/\m M$ over $A/\m$.

\ecoro

\bproof

    If $x_i$ generate $M$ then $\bar x_i$ clearly generate $M$.
    Conversely, let $N=(x_1,\dots,x_n)$ be the submodule generated by $M$, and consider the composition $f\colon N\to M\to M/\m M$.
    $f$'s image is the span of $(\bar x_i)_i$, so all of $M/\m$.
    Thus $f$ is surjective, so $M=N+\m M$, since for every $m\in M$ there is an $n\in N$ such that $f(n)-m\in\m M$ by surjectivity,
    and by above since $J(A)=\m$ this means $M=N$ as desired.
    \qed

\eproof

\bdefn

    Recall the definition of an exact sequence:
    $$ M \xtos fN\xtos gP $$
    is {\emphcolor exact} at $N$ if $\im f=\ker g$.

\edefn

